\documentclass[a4paper,12pt]{article}

\usepackage[french]{babel}
\usepackage[T1]{fontenc}
\usepackage{amsmath,amssymb}
\usepackage[a4paper,margin=1.5cm]{geometry}
\usepackage{multicol}


\usepackage{pgfplots}
\pgfplotsset{compat=newest}


\title{Fiche -- Fonctions de $\mathbb{R}^n$ dans $\mathbb{R}$}
\author{Raphaël Jamann}
\date{}

\begin{document}

    \maketitle
    
    \begin{multicols}{2}[
    \section*{Cadre}
    ]
        On considère une fonction
        $$
        f : \mathbb{R}^n \to \mathbb{R}
        $$

        \begin{itemize}
            \item Entrée : vecteur de $\mathbb{R}^n$
            \item Sortie : un réel
        \end{itemize}

        On note souvent
        $$
        x=(x_1,\dots,x_n)
        $$

        et
        $$
        h=(h_1,\dots,h_n)
        $$

        avec $h$ un vecteur d'incrément.

    \section*{1. Dérivées partielles}

        La dérivée partielle par rapport à $x_i$ est :

        $$
        \frac{\partial f}{\partial x_i}(x)
        =
        \lim_{t\to 0}
        \frac{f(x_1,\dots,x_i+t,\dots,x_n)-f(x)}{t}
        $$

        C'est une dérivée directionnelle particulière, dans la direction de l'axe $x_i$. 
        Par exemple $ \frac{\partial f}{\partial x}(x_0, y_0) = D_{(1,0)}f(x_0, y_0) $

        \subsection*{Nature}

        $$
        \frac{\partial f}{\partial x_i}(x)\in\mathbb{R}
        $$

        C'est donc :
        \begin{itemize}
            \item un scalaire
            \item un nombre réel
        \end{itemize}

        \subsection*{Interprétation}

        Variation de $f$ quand seule la variable $x_i$ varie.

    \section*{2. Gradient}

        Le gradient de $f$ est le vecteur :

        $$
        \nabla f(x)
        =
        \left(
        \frac{\partial f}{\partial x_1}(x),
        \dots,
        \frac{\partial f}{\partial x_n}(x)
        \right)
        $$

        \subsection*{Nature}

        $$
        \nabla f(x)\in\mathbb{R}^n
        $$

        C'est donc :
        \begin{itemize}
            \item un vecteur
        \end{itemize}

        \subsection*{Interprétation}

        Le gradient donne :
        \begin{itemize}
            \item la direction de plus forte croissance
            \item la vitesse maximale de variation
        \end{itemize}

    \section*{3. Dérivée directionnelle}

        Soit $u\in\mathbb{R}^n$ une direction.

        La dérivée directionnelle de $f$ en $x$ selon $u$ est :

        $$
        D_u f(x)
        =
        \lim_{t\to 0}
        \frac{f(x+tu)-f(x)}{t}
        $$

        Si $f$ est différentiable :

        $$
        D_u f(x)=\nabla f(x)\cdot u
        $$

        \subsection*{Nature}

        $$
        D_u f(x)\in\mathbb{R}
        $$

        C'est donc :
        \begin{itemize}
            \item un scalaire
        \end{itemize}

        \subsection*{Interprétation}

        Taux de variation de $f$ dans la direction $u$.

    \section*{4. Différentielle}

        La différentielle de $f$ en $x$ est l'application linéaire :

        $$
        df(x):\mathbb{R}^n\to\mathbb{R}
        $$

        définie par :

        $$
        df(x)(h)
        =
        \sum_{i=1}^n
        \frac{\partial f}{\partial x_i}(x)h_i
        $$

        ou encore :

        $$
        df(x)(h)=\nabla f(x)\cdot h
        $$

        \subsection*{Nature}

        $$
        df(x)\in\mathcal{L}(\mathbb{R}^n,\mathbb{R})
        $$

        C'est donc :
        \begin{itemize}
            \item une application linéaire
        \end{itemize}

        Et :

        $$
        df(x)(h)\in\mathbb{R}
        $$

        est un scalaire.

        \subsection*{Approximation locale}

        Si $f$ est différentiable :

        $$
        f(x+h)
        =
        f(x)+df(x)(h)+o(\|h\|)
        $$

    \section*{5. Hessienne}

        La matrice hessienne est :

        $$
        H_f(x)
        =
        \left(
        \frac{\partial^2 f}{\partial x_i\partial x_j}(x)
        \right)_{1\le i,j\le n}
        $$

        \subsection*{Nature}

        $$
        H_f(x)\in\mathcal{M}_n(\mathbb{R})
        $$

        C'est donc :
        \begin{itemize}
            \item une matrice carrée
        \end{itemize}

    \section*{6. Développement limité d'ordre 2}

        Si $f$ est de classe $C^2$ :

        $$
        f(x+h)
        =
        f(x)
        +
        df(x)(h)
        +
        \frac12 [h]^T H_f(x) [h]
        +
        o(\|h\|^2)
        $$

        ou encore :

        $$
        f(x+h)
        =
        f(x)
        +
        \nabla f(x)\cdot h
        +
        \frac12 [h]^T H_f(x) [h]
        +
        o(\|h\|^2)
        $$

    \end{multicols}

    \section{Extrema locaux et globaux d'une fonction de $\mathbb{R}^n$ dans $\mathbb{R}$}
        \subsection{Recherche points critiques}
            $$ (x, y) \text{ est un point critique} \Longleftrightarrow \overrightarrow{\text{grad}}f(x, y) = \vec{0} \Longleftrightarrow \begin{cases}
                \frac{\partial f}{\partial x}(x, y) = 0 \\
                \frac{\partial f}{\partial y}(x, y) = 0
            \end{cases}
            $$
        
        \subsection{Nature des points critiques}
            La matrice hessienne $H_f(x, y)$ est définie par :
            $$ H_f(x, y) = \begin{pmatrix}
                \frac{\partial^2 f}{\partial x^2}(x, y) & \frac{\partial^2 f}{\partial x \partial y}(x, y) \\
                \frac{\partial^2 f}{\partial y \partial x}(x, y) & \frac{\partial^2 f}{\partial y^2}(x, y)
            \end{pmatrix} $$

            \begin{itemize}
                \item Si $\text{det}(H_f(x, y)) > 0$ et $\text{tr}(H_f(x, y)) > 0$, alors $(x, y)$ est un minimum local.
                \item Si $\text{det}(H_f(x, y)) > 0$ et $\text{tr}(H_f(x, y)) < 0$, alors $(x, y)$ est un maximum local.
                \item Si $\text{det}(H_f(x, y)) < 0$, alors $(x, y)$ est un point selle.
                \item Si $\text{det}(H_f(x, y)) = 0$, $(x, y)$ est dégénéré, il faut utiliser d'autres méthodes pour déterminer la nature du point critique.
            \end{itemize}
            


    \section*{Résumé rapide}

        \begin{center}
        \begin{tabular}{|c|c|}
        \hline
        Objet & Nature \\
        \hline
        $\dfrac{\partial f}{\partial x_i}(x)$ & Scalaire \\
        \hline
        $\nabla f(x)$ & Vecteur de $\mathbb{R}^n$ \\
        \hline
        $D_u f(x)$ & Scalaire \\
        \hline
        $df(x)$ & Application linéaire \\
        \hline
        $df(x)(h)$ & Scalaire \\
        \hline
        $H_f(x)$ & Matrice \\
        \hline
        $h^T H_f(x)h$ & Scalaire \\
        \hline
        \end{tabular}
        \end{center}

\begin{tikzpicture}
  \begin{axis}[samples=30,ticks=none,xlabel={$x$},ylabel={$y$},zlabel={$z$}, title={Point Selle}, clip=false]

    \addplot3[surf,color=blue,opacity=0.5,domain=-2:2,faceted color=black] {x^2-y^2};
    \draw[red, thick, <-] (axis cs:0,0,0) -- (axis cs:1,1,4) node[red, above] {Point selle};
  \end{axis}
\end{tikzpicture}
\quad
\begin{tikzpicture}
  \begin{axis}[samples=40,ticks=none,xlabel={$x$},ylabel={$y$},zlabel={$z$}, title={Extrema Locaux}, clip=false]
    \addplot3[surf,color=blue,opacity=0.5,domain=-2:2,faceted color=black] {4*x*exp(-x^2-y^2)};
    \draw[red, thick, <-] (axis cs:0.707,0,1.7) -- (axis cs:1.5,1,3) node[red, above] {Max local};
    \draw[red, thick, <-] (axis cs:-0.707,0,-1.7) -- (axis cs:-1.5,-1,-3) node[red, below] {Min local};
  \end{axis}
\end{tikzpicture}

\end{document}
